\section{How the assigned repairs will be compared}

\label{sec:repair-study}

\subsection{Eligibility and assignment}

The funded design targets 80 signed findings from at least 20 upstream families, with at most four findings per family. The two planned configurations are Claude Code with Opus 5.5 and the Codex command-line interface with GPT-5.4-mini, both at medium effort. Each configuration receives two standard A assignments per finding, giving 320 standard sessions. Forty paired findings across at least twelve families would add two B assignments per configuration, for 160 further sessions and 480 total. A separate exploratory option would use twelve pairs across at least six families and 368 sessions. It would require another decision. The funded counts describe a target beyond the qualification achieved in Table~\ref{tab:qualification}.

For each paired finding and configuration, we will assign two of four session slots at random to A and two to B. The assignment changes only the example. Each session starts on the same pinned tree with the same report text, tools, environment and budget within that block. The grader will assess A, B and the common hidden H on patches from either arm. Both arms keep the positive assessments private, including the positive check for A when the agent receives B. Before release, we will freeze the signed roster, subsets, evidence versions, choice of examples, software versions, assignment method, run order, limits, grading and retry rules. We will use fresh workspaces, interleave the session order, and check that neither example is available elsewhere under another name in the common environment. The operator receives an assigned brief. The patch grader sees an opaque identifier and the relevant assessment version. Grades and unresolved records will be locked before the assignments are joined to the results.

\subsection{RQ1 and RQ2 on the same attempts}

For RQ1, each configuration's denominator includes every assigned standard A session. We will average its two outcomes within each finding and give the findings equal weight. We will report the numbers assigned, submitting a patch, building, preserving ordinary use and demonstrating the reported correction. Comparing configurations on the same findings allows a reader to see whether both encounter difficulty on the same cases. We will also report a family-weighted rate to show how results change when findings from a heavily represented family carry less weight. A refusal, empty submission, timeout, failed build or demonstrated loss of ordinary use earns no working repair. A valid observation of a persisting violation or a missing required outcome is behavioural failure. If the operator stops a run, we will keep its assigned identifier, original record, stopping reason and any submitted patch. A stop with no submitted patch earns no demonstrated repair. A documented infrastructure retry must follow the rule fixed before release. We will preserve its relationship to the original attempt and never select the better result.

A broken assessment, unknown reach, or contradictory observations leave the affected grade unresolved. A contradiction takes precedence over an automatic verdict. Adjudicators blind to assignment may use existing evidence to apply the frozen rule. If uncertainty remains, the grade stays in the denominator with no demonstrated success. We will report bounds treating unresolved grades first as failures and then as successes, alongside the number and source of the unknowns. A demonstrated build failure, for instance, is a known unsuccessful attempt. Treating it as an unknown would exaggerate the range of rates still consistent with the records.

For RQ2 we will freeze the subset with a separate signed H before the first repair. On precisely the same standard A attempts within this subset, we will report the rate of working repair, the rate also satisfying H, and their difference. We will give counts, denominators, unresolved grades, exclusions, families, relationship types and coverage of the primary roster. There is no applicable H grade outside the subset. The subset has no fixed size and can include findings ineligible for paired assignment. We will repeat the analysis after omitting across-instance observations. A change in rate on the signed subset measures the cost in demonstrated successes of requiring one further behaviour.

\subsection{RQ3 under blind grading}

Each behavioural success indicator requires the patch to build, preserve specified ordinary use, and satisfy the assessment of that behaviour, including its positive evidence. We will grade each patch with the same criteria whichever example the agent received. For each paired finding and configuration, we will compare the proportion repairing A when shown A with the proportion repairing A when shown B. For B we subtract the proportion repairing B when shown A from the proportion repairing B when shown B. We will report the two differences and their equally weighted mean, giving every finding-by-configuration block equal weight. Equivalently, score each patch as half its A success indicator minus half its B indicator. Subtract the mean score in the B arm from the mean in the A arm within each block. A behaviour that succeeds equally often under either assignment contributes zero in its direction, however easy that behaviour is to repair.

We will test the sharp null of no assignment effect with a two-sided randomization test. Each permutation independently preserves the two-versus-two allocation within each block. If sampling permutations, we will retain the observed allocation in the correction and report the seed and Monte Carlo uncertainty. For the primary 95\% interval, in each block we will sum the two within-arm sample variances of patch scores, each divided by its arm size. Summing over blocks and dividing by the squared number of blocks gives the estimated variance of the mean contrast. We will label the resulting normal interval as a large-sample approximation.

We will also give a conservative 95\% interval. For K independent block assignments, each block difference lies between minus one and one. To obtain the half-width, divide twice the natural logarithm of 40 by K and take the square root. Clip the interval to the possible range. For 80 blocks the half-width is about 30 percentage points. If the estimated standard error is zero, this bounded interval governs the decision. We will also assign unresolved grades in the least and most favourable directions for the contrast. In each direction an unknown can occur on either side of a subtraction. Filling every unknown with a success would therefore fail to identify the largest possible advantage for the illustrated behaviour. For the funded design, directional support requires a positive estimate, a governing interval excluding zero, a two-sided randomization p-value below 0.05, and a positive lower bound under unresolved grades. We will report whether the conservative interval excludes zero in every case and show the component differences for each configuration. The exploratory alternative would report estimates and uncertainty without a confirmatory decision.
